\documentclass[10pt,a4paper]{amsart}
\usepackage[margin=19mm]{geometry}
\usepackage{amsmath,amssymb,mathtools}
\usepackage[scr=rsfs,cal=euler]{mathalfa}
\usepackage{xfrac}
\usepackage[unicode,hidelinks,bookmarks=false]{hyperref}
\newcommand{\ie}{i.e.\ }
\newcommand{\cf}{cf.\ }
\newcommand{\resp}{respectively\ }
\newcommand{\Subsection}{\subsection}
\providecommand{\nobreakdash}{\nobreak\hbox{-}}
\setlength{\emergencystretch}{2em}
\multlinegap0pt
\usepackage{bm}
\usepackage{bbm}
\usepackage{dsfont}
\usepackage{xspace}
\usepackage{cancel}
\usepackage{mathtools}
\usepackage{amsthm}
\makeatletter
\@ifundefined{propn}{\newtheorem{propn}{Propn}}{}
\@ifundefined{lemma}{\newtheorem{lemma}[propn]{Lemma}}{}
\@ifundefined{thm}{\newtheorem{thm}[propn]{Theorem}}{}
\makeatother
\newcommand{\D}{\mathbb{D}}
\newcommand{\clb}{\mathcal{B}}
\newcommand{\cle}{\mathcal{E}}
\newcommand{\clf}{\mathcal{F}}
\title[On products and partial isometry of Toeplitz operators with ]{On products and partial isometry of Toeplitz operators with operator-valued symbols}
\author{Srijan Sarkar}
\date{Source: arXiv:2407.10609v2 (CC BY 4.0)}
\begin{document}
\maketitle

\noindent\emph{Source and licence.} On products and partial isometry of Toeplitz operators with operator-valued symbols. Version arXiv:2407.10609v2 (CC BY 4.0), distributed under the Creative Commons Attribution 4.0 International licence, as recorded in the arXiv licence metadata for this exact version and shown on the abstract page https://arxiv.org/abs/2407.10609v2. Full rights notice: licenses/arxiv-2407.10609-CC-BY-4.0.txt.\par\smallskip

\noindent\textit{Editorial scope.} Counted unit: the Thm whose source label is \texttt{tridisc} in the article (source0.tex lines 805--821, source label \texttt{tridisc}), delivered together with the article's own complete original proof of it (source0.tex lines 822--975, one \texttt{proof} environment of 154 lines). No mathematics is written, repaired or re-derived: statement and proof are the article's own text line for line. Required context retained uncounted: rule (lines 696--698). Every definition, display and label of those lines is retained unchanged. Omitted from this package and not redistributed: 1--695, 699--804, 976--1889. Nothing inside a retained excerpt is omitted or altered: every retained body is delivered exactly as the article's lines are written, so no omission receipt of any kind accompanies this package. Citations: the retained excerpts carry no citation command at all, so no bibliography entry of the article is transcribed and no reference is redistributed. Reference adaptation: none was needed and none was made; every reference inside the delivered unit resolves inside it, so no unresolved reference or citation is produced. Printed slips kept literal: no printed word, symbol, hypothesis or number of the counted statement or of its proof was altered. Layout and class: the article's own class is \texttt{amsart} with packages including amssymb, amsmath, amsfonts, latexsym, xcolor, bm; the wrapper is the lane's pinned \texttt{amsart} editorial wrapper at 10pt on A4, which restates the theorem-like environments the retained text needs and adds the article's own macro definitions that the retained text needs (\texttt{\textbackslash D}, \texttt{\textbackslash clb}, \texttt{\textbackslash cle}, \texttt{\textbackslash clf}), with extra wrapper packages (bm, bbm, dsfont, xspace, cancel, mathtools). The delivered numbering is the unit's own. Assets: no retained span contains a figure environment, a graphics-inclusion command, a PDF-page inclusion, a table or a TikZ picture; no image file is copied into this package, the package declares no asset dependency and no image file is redistributed with it. Licence: version arXiv:2407.10609v2 of this article is distributed under the Creative Commons Attribution 4.0 International licence, as recorded in the arXiv licence metadata for this exact version and shown on the abstract page https://arxiv.org/abs/2407.10609v2. A full rights notice is delivered at licenses/arxiv-2407.10609-CC-BY-4.0.txt. Characters: every retained excerpt is pure ASCII, so the delivered engine, XeLaTeX with its default Latin Modern text font, reports no missing character.
\begin{lemma}\label{rule}
Let $\clf, \cle$ be Hilbert spaces and $\Gamma(\bm{z}), \Psi(\bm{z})$ be $\clb(\clf,\cle)$-valued bounded analytic functions on $\D^n$.  Furthermore, let $A,B \subseteq \{1,\ldots,n\}$ such that $A \cup B = \{1,\ldots,n\}$. Then $M_{\Gamma_A} M_{\Psi_B}^*$ is always a Toeplitz operator on $H_{\cle}^2(\D^n)$.
\end{lemma}

\begin{thm}\label{tridisc}
Let $\clf, \cle$ be Hilbert spaces and $\Gamma(\bm{z}), \Psi(\bm{z})$ be $\clb(\clf,\cle)$-valued bounded analytic functions on $\D^3$. Then the following are equivalent \\ \vspace{1mm}
$(i)$~$M_{\Gamma}M_{\Psi}^*$ is a Toeplitz operator on $H_{\cle}^2(\D^3)$, \\ \vspace{1mm}
$(ii)$~$\big( \Gamma(\bm{\lambda}) - \Gamma_k(\bm{\lambda})\big) \big(  \Psi(\bm{\mu}) - \Psi_k(\bm{\mu}) \big)^* = 0$ for all $\bm{\lambda}, \bm{\mu} \in \D^3$ and $k=1,2,3$,\\ \vspace{1mm}
$(iii)$~$M_{\Gamma}M_{\Psi}^*$ admits the following decomposition, \vspace{1mm}
\begin{align*}
M_{\Gamma}M_{\Psi}^*  &=  M_{\Gamma} \Psi(\bm{0})^*\\
&-M_{\Gamma_1} \Psi(\bm{0})^* + M_{\Gamma_1} M_{\Psi_{2,3}}^*\\
&-M_{\Gamma_2} \Psi(\bm{0})^* + M_{\Gamma_2} M_{\Psi_{1,3}}^*\\
&-M_{\Gamma_3} \Psi(\bm{0})^* + M_{\Gamma_3} M_{\Psi_{1,2}}^*\\
&+M_{\Gamma_{1,2}} \Psi(\bm{0})^*  - M_{\Gamma_{1,2}} M_{\Psi_{1,3}}^* -  M_{\Gamma_{1,2}} M_{\Psi_{2,3}}^* + M_{\Gamma_{1,2}} M_{\Psi_{3}}^*\\
&+M_{\Gamma_{1,3}} \Psi(\bm{0})^*  - M_{\Gamma_{1,3}} M_{\Psi_{1,2}}^* -  M_{\Gamma_{1,3}} M_{\Psi_{2,3}}^* + M_{\Gamma_{1,3}} M_{\Psi_{2}}^*\\
&+M_{\Gamma_{2,3}} \Psi(\bm{0})^*  - M_{\Gamma_{2,3}} M_{\Psi_{1,2}}^* - M_{\Gamma_{2,3}} M_{\Psi_{1,3}}^* + M_{\Gamma_{2,3}} M_{\Psi_1}^* \\
&+\Gamma(\bm{0}) M_{\Psi}^*  - \Gamma(\bm{0}) M_{\Psi_{1}}^* -  \Gamma(\bm{0}) M_{\Psi_{2}}^* - \Gamma(\bm{0}) M_{\Psi_{3}}^*\\
&+\Gamma(\bm{0}) M_{\Psi_{1,2}}^*  - \Gamma(\bm{0}) M_{\Psi_{1,3}}^* -  \Gamma(\bm{0}) M_{\Psi_{2,3}}^* - \Gamma(\bm{0}) \Psi(\bm{0})^*.
\end{align*}
\end{thm}

\begin{proof}
The $(i) \iff (ii)$ direction is the same as in the previous theorem.  For the direction $(ii) \implies (iii)$, let us start with condition $(ii)$, when $k=1$, that is,
\[
\big( \Gamma(\bm{\lambda}) - \Gamma_1(\bm{\lambda}) \big) \big( \Psi(\bm{\mu}) - \Psi_1(\bm{\mu}) \big)^* = 0.
\]
This gives the following decomposition.
\begin{equation}\label{da1}
\Gamma(\bm{\lambda})  \Psi(\bm{\mu})^* =  \Gamma(\bm{\lambda}) \Psi_1(\bm{\mu})^* + \Gamma_1(\bm{\lambda})\Psi(\bm{\mu})^* - \Gamma_1(\bm{\lambda}) \Psi_1(\bm{\mu})^*.
\end{equation}
Again using condition $(ii)$ when $k=2$ and $\mu_1 = 0$, we get
\[
\big( \Gamma(\bm{\lambda}) - \Gamma_2(\bm{\lambda}) \big) \big( \Psi_{1}(\bm{\mu}) - \Psi_{1,2}(\bm{\mu}) \big)^* = 0.
\]
This implies that,
\[
 \Gamma(\bm{\lambda}) \Psi_1(\bm{\mu})^* =   \Gamma(\bm{\lambda}) \Psi_{1,2}(\bm{\mu})^* + \Gamma_2(\bm{\lambda}) \Psi_1(\bm{\mu})^* - \Gamma_2(\bm{\lambda}) \Psi_{1,2}(\bm{\mu})^* 
\]
Since the above condition is assumed to be true for all $\bm{\lambda}, \bm{\mu}$ in $\D^3$, we can put $\lambda_1 = 0$ to get
\[
\Gamma_1(\bm{\lambda}) \Psi_{1}(\bm{\mu})^* = \Gamma_1(\bm{\lambda}) \Psi_{1,2}(\bm{\mu})^*
+ \Gamma_{1,2}(\bm{\lambda})  \Psi_{1}(\bm{\mu})^* - \Gamma_{1,2}(\bm{\lambda}) \Psi_{1,2}(\bm{\mu})^*.
\]
Using condition $(ii)$, when $k=2$ and $\lambda_1 = 0$ gives
\[
\big( \Gamma_1(\bm{\lambda}) - \Gamma_{1,2}(\bm{\lambda}) \big) \big( \Psi(\bm{\mu}) - \Psi_{2}(\bm{\mu}) \big)^* = 0.
\]
\begin{equation*}
\Gamma_1(\bm{\lambda})\Psi(\bm{\mu})^* =  \Gamma_1(\bm{\lambda})\Psi_2(\bm{\mu})^* + \Gamma_{1,2}(\bm{\lambda}) \Psi(\bm{\mu})^* - \Gamma_{1,2}(\bm{\lambda}) \Psi_{2}(\bm{\mu})^*.
\end{equation*}
Incorporating all the above three decompositions inside condition (\ref{da1}), we get
\begin{equation}\label{da2}
\begin{split}
\Gamma(\bm{\lambda})  \Psi(\bm{\mu})^* &=  \Gamma(\bm{\lambda}) \Psi_{1,2}(\bm{\mu})^* + \Gamma_2(\bm{\lambda}) \Psi_1(\bm{\mu})^* - \Gamma_2(\bm{\lambda}) \Psi_{1,2}(\bm{\mu})^* \\
&+ \Gamma_1(\bm{\lambda})\Psi_2(\bm{\mu})^* + \Gamma_{1,2}(\bm{\lambda}) \Psi(\bm{\mu})^* - \Gamma_{1,2}(\bm{\lambda}) \Psi_{2}(\bm{\mu})^* \\
&- \Gamma_1(\bm{\lambda}) \Psi_{1,2}(\bm{\mu})^*
- \Gamma_{1,2}(\bm{\lambda})  \Psi_{1}(\bm{\mu})^* + \Gamma_{1,2}(\bm{\lambda}) \Psi_{1,2}(\bm{\mu})^*.
\end{split}
\end{equation}
Using condition $(ii)$, when $k=3$ and $\mu_1=\mu_2=0$, we get
\[
\big( \Gamma(\bm{\lambda}) - \Gamma_3(\bm{\lambda}) \big) \big( \Psi_{1,2}(\bm{\mu}) - \Psi (\bm{0}) \big)^* = 0.
\]
This gives the following decomposition.
\[
 \Gamma(\bm{\lambda})\Psi_{1,2}(\bm{\mu})^* =    \Gamma_3(\bm{\lambda}) \Psi_{1,2}(\bm{\mu})^* + \Gamma(\bm{\lambda}) \Psi(\bm{0})^* - \Gamma_3(\bm{\lambda}) \Psi(\bm{0})^*.
\]
Again using condition $(ii)$, when $k=3$ and $\lambda_1 = \lambda_2 = 0$, we get
\[
\big( \Gamma_{1,2}(\bm{\lambda}) - \Gamma (\bm{0}) \big) \big( \Psi(\bm{\mu}) - \Psi_3(\bm{\mu}) \big)^* = 0.
\]
This gives the following decomposition.
\[
\Gamma_{1,2}(\bm{\lambda}) \Psi(\bm{\mu})^* = \Gamma_{1,2}(\bm{\lambda})\Psi_{3}(\bm{\mu})^* + \Gamma(\bm{0})\Psi(\bm{\mu})^* - \Gamma(\bm{0})  \Psi_{3}(\bm{\mu})^*.
\]
Incorporating the above decompositions in condition (\ref{da2}), we get
\begin{equation}\label{da3}
\begin{split}
\Gamma(\bm{\lambda})  \Psi(\bm{\mu})^* &=  \Gamma_3(\bm{\lambda}) \Psi_{1,2}(\bm{\mu})^* + \Gamma(\bm{\lambda}) \Psi(\bm{0})^* - \Gamma_3(\bm{\lambda}) \Psi(\bm{0})^* \\
&+ \Gamma_2(\bm{\lambda}) \Psi_1(\bm{\mu})^* - \Gamma_2(\bm{\lambda}) \Psi_{1,2}(\bm{\mu})^* \\
&+ \Gamma_1(\bm{\lambda})\Psi_2(\bm{\mu})^*\\ 
&+ \Gamma_{1,2}(\bm{\lambda})\Psi_{3}(\bm{\mu})^* + \Gamma(\bm{0})\Psi(\bm{\mu})^* - \Gamma(\bm{0})  \Psi_{3}(\bm{\mu})^*\\
&- \Gamma_{1,2}(\bm{\lambda}) \Psi_{2}(\bm{\mu})^* \\
&- \Gamma_1(\bm{\lambda}) \Psi_{1,2}(\bm{\mu})^*
- \Gamma_{1,2}(\bm{\lambda})  \Psi_{1}(\bm{\mu})^* + \Gamma_{1,2}(\bm{\lambda}) \Psi_{1,2}(\bm{\mu})^*.
\end{split}
\end{equation}
Condition $(ii)$ when $k=3$, and $\lambda_1 = \mu_2 = 0$, gives us
\[
\big( \Gamma_{1}(\bm{\lambda}) - \Gamma_{1,3}(\bm{\lambda}) \big) \big( \Psi_2(\bm{\mu}) - \Psi_{2,3}(\bm{\mu}) \big)^* = 0,
\]
which provides the following decomposition.
\[
\Gamma_1(\bm{\lambda}) \Psi_{2}(\bm{\mu})^* = \Gamma_1(\bm{\lambda})  \Psi_{2,3}(\bm{\mu})^* +  \Gamma_{1,3}(\bm{\lambda}) \Psi_{2}(\bm{\mu})^* - \Gamma_{1,3}(\bm{\lambda}) \Psi_{2,3}(\bm{\mu})^*.
\]
Now, we will separately consider the following cases: $(a)~\mu_1 = 0 $, $(b)~\lambda_2 = 0$ and $(c)~\lambda_2 = \mu_1 = 0$, in the above condition to get the following three decompositions.
\[
\Gamma_1(\bm{\lambda})\Psi_{1,2}(\bm{\mu})^* =  \Gamma_1(\bm{\lambda}) \Psi(\bm{0})^* +  \Gamma_{1,3}(\bm{\lambda}) \Psi_{1,2}(\bm{\mu})^* - \Gamma_{1,3}(\bm{\lambda}) \Psi(\bm{0})^*. \tag{\text{a}}
\]
\[
\Gamma_{1,2}(\bm{\lambda}) \Psi_{2}(\bm{\mu})^* =   \Gamma_{1,2}(\bm{\lambda}) \Psi_{2,3}(\bm{\mu})^* + \Gamma(\bm{0}) \Psi_{2}(\bm{\mu})^* - \Gamma(\bm{0}) \Psi_{2,3}(\bm{\mu})^*. \tag{\text{b}}
\]
\[
\Gamma_{1,2}(\bm{\lambda})  \Psi_{1,2}(\bm{\mu})^* = \Gamma_{1,2}(\bm{\lambda}) \Psi(\bm{0})^* + \Gamma(\bm{0}) \Psi_{1,2}(\bm{\mu})^* - \Gamma(\bm{0}) \Psi(\bm{0})^*. \tag{\text{c}}
\]
Incorporating the above identities in condition (\ref{da3}), we get
\begin{equation}\label{da4}
\begin{split}
\Gamma(\bm{\lambda})  \Psi(\bm{\mu})^* &=  \Gamma_3(\bm{\lambda}) \Psi_{1,2}(\bm{\mu})^* + \Gamma(\bm{\lambda}) \Psi(\bm{0})^* - \Gamma_3(\bm{\lambda}) \Psi(\bm{0})^* \\
&+ \Gamma_2(\bm{\lambda}) \Psi_1(\bm{\mu})^* - \Gamma_2(\bm{\lambda}) \Psi_{1,2}(\bm{\mu})^* \\
&+ \Gamma_1(\bm{\lambda})  \Psi_{2,3}(\bm{\mu})^* +  \Gamma_{1,3}(\bm{\lambda}) \Psi_{2}(\bm{\mu})^* - \Gamma_{1,3}(\bm{\lambda}) \Psi_{2,3}(\bm{\mu})^* \\
&+ \Gamma_{1,2}(\bm{\lambda})\Psi_{3}(\bm{\mu})^* + \Gamma(\bm{0})\Psi(\bm{\mu})^* - \Gamma(\bm{0})  \Psi_{3}(\bm{\mu})^*\\
&- \Gamma_{1,2}(\bm{\lambda}) \Psi_{2,3}(\bm{\mu})^* - \Gamma(\bm{0}) \Psi_{2}(\bm{\mu})^* + \Gamma(\bm{0}) \Psi_{2,3}(\bm{\mu})^* \\
&- \Gamma_1(\bm{\lambda}) \Psi(\bm{0})^* - \Gamma_{1,3}(\bm{\lambda}) \Psi_{1,2}(\bm{\mu})^* + \Gamma_{1,3}(\bm{\lambda}) \Psi(\bm{0})^*.\\
&- \Gamma_{1,2}(\bm{\lambda})  \Psi_{1}(\bm{\mu})^* \\
&+ \Gamma_{1,2}(\bm{\lambda}) \Psi(\bm{0})^* + \Gamma(\bm{0}) \Psi_{1,2}(\bm{\mu})^* - \Gamma(\bm{0}) \Psi(\bm{0})^*.
\end{split}
\end{equation}
In a similar manner, using condition $(ii)$, when $k=3$ and $\lambda_2 = \mu_1 = 0$ gives us 
\[
\big( \Gamma_{2}(\bm{\lambda}) - \Gamma_{2,3}(\bm{\lambda}) \big) \big( \Psi_1(\bm{\mu}) - \Psi_{1,3}(\bm{\mu}) \big)^* = 0,
\]
and hence,  the following decomposition.
\[
\Gamma_2(\bm{\lambda}) \Psi_{1}(\bm{\mu})^* =  \Gamma_2(\bm{\lambda})\Psi_{1,3}(\bm{\mu})^* + \Gamma_{2,3}(\bm{\lambda}) \Psi_{1}(\bm{\mu})^* - \Gamma_{2,3}(\bm{\lambda}) \Psi_{1,3}(\bm{\mu})^*. 
\]
Considering the following cases separately: $(d)~\lambda_1 =0 $, $(e)~\mu_2 = 0$ in the above condition gives us the following two decompositions.
\[
\Gamma_{1,2}(\bm{\lambda}) \Psi_{1}(\bm{\mu})^* =  \Gamma_{1,2}(\bm{\lambda})\Psi_{1,3}(\bm{\mu})^* + \Gamma (\bm{0}) \Psi_{1}(\bm{\mu})^* - \Gamma (\bm{0}) \Psi_{1,3}(\bm{\mu})^*. \tag{\text{d}}
\]
\[
\Gamma_2(\bm{\lambda}) \Psi_{1,2}(\bm{\mu})^* =  \Gamma_2(\bm{\lambda}) \Psi(\bm{0})^*+ \Gamma_{2,3}(\bm{\lambda}) \Psi_{1,2}(\bm{\mu})^* - \Gamma_{2,3}(\bm{\lambda}) \Psi(\bm{0})^*.  \tag{\text{e}}
\]
Incorporating the above decompositions inside condition (\ref{da4}), gives us
\begin{equation}\label{da5}
\begin{split}
\Gamma(\bm{\lambda})  \Psi(\bm{\mu})^* &=  \Gamma_3(\bm{\lambda}) \Psi_{1,2}(\bm{\mu})^* + \Gamma(\bm{\lambda}) \Psi(\bm{0})^* - \Gamma_3(\bm{\lambda}) \Psi(\bm{0})^* \\
&+ \Gamma_2(\bm{\lambda})\Psi_{1,3}(\bm{\mu})^* + \Gamma_{2,3}(\bm{\lambda}) \Psi_{1}(\bm{\mu})^* - \Gamma_{2,3}(\bm{\lambda}) \Psi_{1,3}(\bm{\mu})^*. \\
&- \Gamma_2(\bm{\lambda}) \Psi(\bm{0})^*- \Gamma_{2,3}(\bm{\lambda}) \Psi_{1,2}(\bm{\mu})^* + \Gamma_{2,3}(\bm{\lambda}) \Psi(\bm{0})^*\\
&+ \Gamma_1(\bm{\lambda})  \Psi_{2,3}(\bm{\mu})^* +  \Gamma_{1,3}(\bm{\lambda}) \Psi_{2}(\bm{\mu})^* - \Gamma_{1,3}(\bm{\lambda}) \Psi_{2,3}(\bm{\mu})^* \\
&+ \Gamma_{1,2}(\bm{\lambda})\Psi_{3}(\bm{\mu})^* + \Gamma(\bm{0})\Psi(\bm{\mu})^* - \Gamma(\bm{0})  \Psi_{3}(\bm{\mu})^*\\
\
&- \Gamma_{1,2}(\bm{\lambda}) \Psi_{2,3}(\bm{\mu})^* - \Gamma(\bm{0}) \Psi_{2}(\bm{\mu})^* + \Gamma(\bm{0}) \Psi_{2,3}(\bm{\mu})^* \\
&- \Gamma_1(\bm{\lambda}) \Psi(\bm{0})^* - \Gamma_{1,3}(\bm{\lambda}) \Psi_{1,2}(\bm{\mu})^* + \Gamma_{1,3}(\bm{\lambda}) \Psi(\bm{0})^*.\\
&- \Gamma_{1,2}(\bm{\lambda})\Psi_{1,3}(\bm{\mu})^* - \Gamma (\bm{0}) \Psi_{1}(\bm{\mu})^* + \Gamma (\bm{0}) \Psi_{1,3}(\bm{\mu})^*. \\
&+ \Gamma_{1,2}(\bm{\lambda}) \Psi(\bm{0})^* + \Gamma(\bm{0}) \Psi_{1,2}(\bm{\mu})^* - \Gamma(\bm{0}) \Psi(\bm{0})^*.
\end{split}
\end{equation}
Therefore, for any $\eta, \zeta \in \cle$, we get
\begin{align*}
\langle M_{\Gamma}M_{\Psi}^* k_{\bm{\mu}} \eta, k_{\bm{\lambda}} \zeta \rangle &= k_{\bm{{\mu}}}(\bm{\lambda}) \langle  \Gamma(\bm{\lambda})\Psi(\bm{\mu})^*  \eta,  \zeta \rangle \\
&= k_{\bm{{\mu}}}(\bm{\lambda}) \Big( \langle  \Gamma_3(\bm{\lambda}) \Psi_{1,2}(\bm{\mu})^* \eta,  \zeta \rangle + \langle  \Gamma(\bm{\lambda}) \Psi(\bm{0})^* \eta,  \zeta \rangle - \langle  \Gamma_3(\bm{\lambda}) \Psi(\bm{0})^* \eta,  \zeta \rangle \\
&+ \langle  \Gamma_2(\bm{\lambda})\Psi_{1,3}(\bm{\mu})^* \eta,  \zeta \rangle + \langle  \Gamma_{2,3}(\bm{\lambda}) \Psi_{1}(\bm{\mu})^* \eta,  \zeta \rangle - \langle \Gamma_{2,3}(\bm{\lambda}) \Psi_{1,3}(\bm{\mu})^* \eta,  \zeta \rangle. \\
&- \langle \Gamma_2(\bm{\lambda}) \Psi(\bm{0})^* \eta,  \zeta \rangle - \langle \Gamma_{2,3}(\bm{\lambda}) \Psi_{1,2}(\bm{\mu})^* \eta,  \zeta \rangle + \langle \Gamma_{2,3}(\bm{\lambda}) \Psi(\bm{0})^* \eta,  \zeta \rangle. \\
&+ \langle \Gamma_1(\bm{\lambda}) \Psi_{2,3}(\bm{\mu})^* \eta,  \zeta \rangle +  \langle \Gamma_{1,3}(\bm{\lambda}) \Psi_{2}(\bm{\mu})^* \eta,  \zeta \rangle - \langle \Gamma_{1,3}(\bm{\lambda}) \Psi_{2,3}(\bm{\mu})^* \eta,  \zeta \rangle \\
&+ \langle \Gamma_{1,2}(\bm{\lambda})\Psi_{3}(\bm{\mu})^* \eta,  \zeta \rangle + \langle \Gamma(\bm{0})\Psi(\bm{\mu})^* \eta,  \zeta \rangle - \langle \Gamma(\bm{0})  \Psi_{3}(\bm{\mu})^* \eta,  \zeta \rangle\\
&- \langle \Gamma_{1,2}(\bm{\lambda}) \Psi_{2,3}(\bm{\mu})^* \eta,  \zeta \rangle - \langle \Gamma(\bm{0}) \Psi_{2}(\bm{\mu})^* \eta,  \zeta \rangle + \langle \Gamma(\bm{0}) \Psi_{2,3}(\bm{\mu})^* \eta,  \zeta \rangle \\
&- \langle \Gamma_1(\bm{\lambda}) \Psi(\bm{0})^* \eta,  \zeta \rangle - \langle \Gamma_{1,3}(\bm{\lambda}) \Psi_{1,2}(\bm{\mu})^* \eta,  \zeta \rangle + \langle \Gamma_{1,3}(\bm{\lambda}) \Psi(\bm{0})^* \eta,  \zeta \rangle.\\
&- \langle \Gamma_{1,2}(\bm{\lambda})\Psi_{1,3}(\bm{\mu})^* \eta,  \zeta \rangle - \langle \Gamma (\bm{0}) \Psi_{1}(\bm{\mu})^* \eta,  \zeta \rangle + \langle \Gamma (\bm{0}) \Psi_{1,3}(\bm{\mu})^* \eta,  \zeta \rangle . \\
&+ \langle \Gamma_{1,2}(\bm{\lambda}) \Psi(\bm{0})^* \eta,  \zeta \rangle + \langle \Gamma(\bm{0}) \Psi_{1,2}(\bm{\mu})^* \eta,  \zeta \rangle - \langle \Gamma(\bm{0}) \Psi(\bm{0})^* \eta,  \zeta \rangle \Big).
\end{align*}
Since the above is true for all $\bm{\lambda}, \bm{\mu} \in \D^3$, and $\eta, \zeta \in \cle$, we get
\begin{align*}
M_{\Gamma}M_{\Psi}^*  &=  M_{\Gamma_3}M_{\Psi_{1,2}}^* + M_{\Gamma} \Psi(\bm{0})^* - M_{\Gamma_3} \Psi(\bm{0})^*\\
&+ M_{\Gamma_2} M_{\Psi_{1,3}}^* + M_{\Gamma_{2,3}} M_{\Psi_{1}}^* - M_{\Gamma_{2,3}}  M_{\Psi_{2,3}}^*\\
&-M_{\Gamma_2} \Psi(\bm{0})^* - M_{\Gamma_{2,3}} M_{\Psi_{1,2}}^* + M_{\Gamma_{2,3}}  \Psi(\bm{0})^*\\
&+M_{\Gamma_{1,2}} M_{\Psi_3}^* + \Gamma(\bm{0}) M_{\Psi}^* - \Gamma(\bm{0}) M_{\Psi_3}^*\\
&-M_{\Gamma_{1,2}} M_{\Psi_{2,3}}^*  - \Gamma(\bm{0}) M_{\Psi_{2}}^* +  \Gamma(\bm{0}) M_{\Psi_{2,3}}^*\\
&-M_{\Gamma_{1}} \Psi(\bm{0})^*  - M_{\Gamma_{1,3}} M_{\Psi_{1,2}}^* +  M_{\Gamma_{1,3}} \Psi(\bm{0})^* \\
&-M_{\Gamma_{1,2}} M_{\Psi_{1,3}}^*  - \Gamma(\bm{0}) M_{\Psi_{1}}^* +  \Gamma(\bm{0}) M_{\Psi_{1,3}}^* \\
&+M_{\Gamma_{1,2}} \Psi(\bm{0})^*  + \Gamma(\bm{0}) M_{\Psi_{1,2}}^* -  \Gamma(\bm{0}) \Psi(\bm{0})^*.
\end{align*}
This completes the proof for the direction $(ii) \implies (iii)$. For the direction $(iii) \implies (i)$, it is easy to observe using Lemma \ref{rule} that all the operators on the right-hand side of the above decomposition are Toeplitz operators. Thus, $M_{\Gamma}M_{\Psi}^*$ must be a Toeplitz operator as well.
This completes the proof.
\end{proof}

\end{document}
